\documentclass[11pt]{article}

\usepackage[margin=1.1in]{geometry}
\usepackage{amsmath,amssymb,amsthm}
\usepackage{booktabs}
\usepackage{enumitem}
\usepackage{xspace}
\usepackage{xcolor}
\usepackage[colorlinks=true,linkcolor=blue!50!black,citecolor=blue!50!black,urlcolor=blue!50!black]{hyperref}

% ---------------------------------------------------------------------------
% Macros
% ---------------------------------------------------------------------------
\newcommand{\F}{\mathbb{F}}                       % the extension field F_{2^128}
\newcommand{\Ftwo}{\F_2}
\newcommand{\eq}{\mathsf{eq}}
\newcommand{\nnz}{\mathsf{nnz}}
\newcommand{\ip}[2]{\langle #1, #2 \rangle}
\newcommand{\code}[1]{\texttt{#1}}
\newcommand{\ZC}{\textsf{Zerocheck}\xspace}
\newcommand{\LC}{\textsf{Lincheck}\xspace}
\newcommand{\pre}{\mathrm{pre}}
\newcommand{\out}{\mathrm{out}}
\newcommand{\inn}{\mathrm{in}}

\newtheorem{lemma}{Lemma}
\newtheorem{remark}{Remark}
\theoremstyle{definition}
\newtheorem{definition}{Definition}

\title{Multi-Table Proving in Flock:\\ One Proof for Heterogeneous Hash Batches\\[0.6ex]
\large Design document --- working draft v0.2}
\author{Flock}
\date{July 22, 2026}

\begin{document}
\maketitle

\begin{abstract}
Flock currently produces one proof per hash function: a proof attests to $n$
invocations of \emph{one} of BLAKE3, SHA-256, or Keccak-f[1600]. This document
designs the generalization to a \emph{single} proof attesting to a
heterogeneous batch --- e.g.\ $1000$ BLAKE3 compressions \emph{and} $4096$
SHA-256 compressions, independently invoked, with the counts chosen at prove
time --- in one proof with one commitment and one opening. The design views
the statement as a collection of \emph{tables} (one per hash type; rows =
invocations, columns = the trace of one invocation) composed into one
\emph{block-diagonal R1CS} over a \textbf{static, capacity-sized address
space}: each type owns a fixed slot sized to a maximum invocation count, the
per-proof counts are public inputs, and undeclared rows are zero --- they
satisfy the constraints identically and cost the prover essentially nothing.
The central observation is that running Flock's existing single-instance
zerocheck over the composed instance is \emph{mathematically identical} to a
batched per-table zerocheck with $\eq$-derived batching coefficients --- so
batching comes for free, with plain single-instance soundness and almost no
new zerocheck code. The genuinely new protocol work concentrates in the
lincheck, which needs a small multi-size sumcheck over the per-table column
dimensions. The commitment layer is the jagged (stacked-columns) scheme:
only the actual data is committed, with the static address space existing
purely virtually. There is one protocol mode; fixed workloads are simply the
full-utilization special case.
\end{abstract}

\tableofcontents

% ===========================================================================
\section{Goals and non-goals}\label{sec:goals}

\paragraph{Goal.} One proof system execution that proves a batch containing
invocations of \emph{several different} hash functions, with the number of
invocations of each chosen at prove time. The invocations are
\textbf{independent}: no hash output feeds another hash input across (or
within) types; each row is a self-contained invocation. The multi-table
core deliberately binds \emph{no} per-invocation input/output --- all
dataflow, between tables and to the public statement, is deferred to a
planned lookup/bus layer (\S\ref{sec:statement}, \S\ref{sec:open}).

\paragraph{Design priorities} (decided up front, in order):
\begin{enumerate}[nosep]
  \item \textbf{Prover throughput.} Hashes proven per second is the headline
        metric. Per-hash witness-generation kernels and the NEON/AVX hot
        paths must not be perturbed, and the prover must pay for
        \emph{declared} work, never for capacity.
  \item \textbf{Gradual landing.} Small reviewable steps; the existing
        single-hash path stays green throughout; each phase is independently
        shippable and benchmarkable.
  \item \textbf{Static structure.} The set of hash types, their encoders,
        their widths, and their capacities come from a small predeclared
        \emph{registry}. Everything structural about the verifier --- round
        counts, control flow, precomputed digests --- depends only on the
        registry, never on the per-proof counts, which enter as scalars
        only. (This is what keeps the door open to fixed-circuit recursion
        or aggregation later.)
\end{enumerate}

\paragraph{Non-goals (for now).}
Cross-table dataflow and per-invocation public IO --- both deliberately
excluded from the multi-table core and slated for a subsequent lookup/bus
layer (\S\ref{sec:open}); a dynamic \emph{type} registry (adding a new
hash type remains a registry/code change); changing the \emph{internal}
protocol hash (the SHA-256 used for the Merkle commitment, Fiat--Shamir,
and proof-of-work grinding is untouched by this design --- ``hash'' in this
document always means the hash being \emph{proven}).

\paragraph{One mode.} An earlier draft of this design carried two protocol
modes: a \emph{fixed-shape} mode (counts predeclared, witness committed
directly) and the dynamic mode. The fixed-shape mode is expressively the
special case ``counts $=$ capacities'' of the dynamic mode, and carrying
both meant two commitment layers, two transports, two wire formats, and a
doubled audit surface. This draft specifies \textbf{the dynamic mode only}.
The direct-commitment path survives as a development harness and fallback
(Remark~\ref{rem:harness}), not as a protocol mode.

% ===========================================================================
\section{Preliminaries: the current single-table system}\label{sec:prelim}

\subsection{Notation}

$\F = \F_{2^{128}}$ is the extension field; witnesses are $\Ftwo$-valued
(bits). For $\tau, x \in \F^m$,
\[
  \eq(\tau, x) \;=\; \prod_{i=1}^{m} \bigl(\tau_i x_i + (1-\tau_i)(1-x_i)\bigr),
\]
the multilinear extension of the Kronecker delta on the Boolean cube. We use
throughout the identity
\begin{equation}\label{eq:eqsum}
  \sum_{x \in \{0,1\}^m} \eq(\tau, x) \;=\; 1
  \qquad\text{for every } \tau \in \F^m,
\end{equation}
which holds because each factor sums to $\tau_i + (1-\tau_i) = 1$ over
$x_i \in \{0,1\}$. For a vector $v \in \F^{2^m}$, $\hat v$ denotes its
multilinear extension (MLE); index bits are little-endian (bit $0$ is the
least significant). For a point $r \in \F^m$ and a bit range $[a,b)$ we write
$r[a{:}b]$ for the corresponding subvector of coordinates.

\subsection{The single-table instance}\label{sec:single}

A Flock statement today is a \emph{block R1CS}
(\code{flock-core/src/r1cs.rs}, \code{struct BlockR1cs}): one \emph{base
block} of sparse $\Ftwo$-matrices $(A_0, B_0, C_0)$, each $k \times k$ with
$k = 2^{\kappa}$, with $C_0 = I$ (``circuit form'': the $C$-side just reads
the witness), replicated $n = 2^{\nu}$ times. The full constraint matrix is
implicitly
\[
  A \;=\; I_{n} \otimes A_0, \qquad B \;=\; I_{n} \otimes B_0,
  \qquad C \;=\; I_{n} \otimes C_0 = I,
\]
and the witness $z \in \Ftwo^{2^{m}}$, $m = \nu + \kappa$, satisfies
$(Az) \circ (Bz) = Cz$ (Hadamard product), i.e.\ each block independently
satisfies the base relation. Only the first $u \le k$ columns of a block are
``useful''; columns $[u, k)$ are zero padding, recorded by a single scalar
\code{useful\_bits}.

\paragraph{The table view.} This is exactly a \emph{table}: $n$ rows (one per
hash invocation), $k$ columns (the trace of one invocation --- for the current
encoders, $\kappa = 14$ for BLAKE3, $15$ for SHA-256, $16$ for Keccak, $17$
for the 3-wide Keccak variant; all share the zerocheck skip parameter
$k_{\mathrm{skip}} = 6$). The witness index splits as
\[
  x \;=\; (\underbrace{x[0{:}\kappa]}_{\text{column / ``inner''}},\;
           \underbrace{x[\kappa{:}m]}_{\text{row / ``outer''}}).
\]
This is today's default \emph{row-major} convention (one invocation's
trace is contiguous); the multi-table design adopts the code's other,
chunk-column-major convention (\code{WitnessLayout::BatchMajor},
\S\ref{sec:union}).

\subsection{The PIOP, briefly}\label{sec:pioponetable}

Let $a = Az$, $b = Bz$, $c = Cz = z$ over the $2^m$-point domain. The proof
runs:

\paragraph{\ZC (\code{zerocheck.rs}).} Proves $\hat a(x)\hat b(x) = \hat c(x)$
for all Boolean $x$: the verifier samples $\tau \in \F^{m}$, the prover runs a
sumcheck for
\begin{equation}\label{eq:zc-single}
  \sum_{x \in \{0,1\}^{m}} \eq(\tau, x)\,
     \bigl(\hat a(x)\hat b(x) - \hat c(x)\bigr) \;=\; 0 .
\end{equation}
The low $k_{\mathrm{skip}} = 6$ variables are bound in a single
\emph{univariate-skip} round (one univariate over a $2^{6}$-point domain),
which is where most of the prover's win from $\Ftwo$-valued data lives; the
remaining $m - 6$ rounds are ordinary. The output is a random point
$r \in \F^{m}$ and claims $\hat a(r), \hat b(r), \hat c(r)$. Since $C = I$,
$\hat c(r) = \hat z(r)$ is already a witness evaluation claim. A
\code{PaddingSpec} $= (\kappa, u)$ lets the prover skip work on the uniform
zero-padding columns.

\paragraph{\LC (\code{lincheck.rs}).} Reduces $\hat a(r)$ and $\hat b(r)$
(batched by a fresh $\alpha$: the claim is $\alpha\hat a(r) + \hat b(r)$) to a
witness evaluation. The load-bearing structure is the factorization of the
replicated matrix's MLE:
\begin{equation}\label{eq:factor-single}
  \hat A\bigl((r_{\out}, r_{\inn}), (c_{\out}, c_{\inn})\bigr)
  \;=\; \eq(r_{\out}, c_{\out}) \cdot \hat A_0(r_{\inn}, c_{\inn}),
\end{equation}
so that
\[
  \hat a(r) = \sum_{c} \hat A(r, c)\,\hat z(c)
  = \sum_{c_{\inn} \in \{0,1\}^{\kappa}}
      \hat A_0(r_{\inn}, c_{\inn})\;
      \underbrace{\hat z(c_{\inn}, r_{\out})}_{\text{partial eval.\ of } \hat z},
\]
where the sum over $c_{\out}$ \emph{collapsed} into a partial evaluation of
the witness MLE at the row point $r_{\out}$. The prover materializes
$y(c_{\inn}) := \hat z(c_{\inn}, r_{\out})$ (one $O(2^m)$ fold of the witness)
and the \emph{comb} vector
$\mathrm{comb}(c_{\inn}) = \alpha\,(\eq_{r_{\inn}}^{\top} A_0)(c_{\inn})
 + (\eq_{r_{\inn}}^{\top} B_0)(c_{\inn})$,
then runs a $\kappa$-bit sumcheck for
$\sum_{c_{\inn}} \mathrm{comb}(c_{\inn})\, y(c_{\inn})$. The verifier's matrix
work is $O(\nnz(A_0) + \nnz(B_0))$, \textbf{independent of the row count} ---
it never touches the expanded matrix. (In code the comb evaluation goes
through \code{trait LincheckCircuit}; Keccak implements it with a custom
walker rather than literal sparse matrices. Everything below interacts with
that trait, not with matrices directly.)

\paragraph{Claims, glue, opening.} All evaluation claims on $\hat z$ --- from
\ZC ($\hat c(r)$), from \LC, and from the statement-binding arguments (e.g.\
today's chain argument, which constrains input/output \emph{regions} of the
committed witness via an auxiliary shift-sumcheck; note it is claim-level
``glue'', \emph{not} extra R1CS rows) --- are carried to the committed
polynomial by a two-stage machinery: \emph{ring-switching}
(\code{pcs/ring\_switch.rs}) eliminates the seven in-word packing
coordinates --- including the univariate skip coordinate, which is why
downstream layers never see quirky structure --- turning each bit-level
claim into a multilinear claim on the \emph{packed} ($\F$-word)
polynomial; the packed claims are then batched
(\code{pcs::open\_batch\_mixed\ldots}, label
\code{flock-pcs-open-batch-v0}) into one Ligerito opening of the single
committed witness polynomial.

The design below preserves every element of this pipeline and generalizes the
instance it runs on.

% ===========================================================================
\section{The multi-table model}\label{sec:model}

\begin{definition}[Type registry]\label{def:registry}
The \emph{registry} fixes, once, an ordered list of \emph{table types}
$t = 1, \dots, T$. Type $t$ consists of: a base block
$(A_t, B_t, C_t{=}I)$ of width $k_t = 2^{\kappa_t}$ with $u_t \le k_t$
useful columns, and a \emph{row capacity} $2^{\nu_t^{\max}}$ --- the padded
maximum: the largest number of invocations of type $t$ a single proof may
declare, sizing the type's (virtual) slot once and for all. The number of
rows \emph{actually used} is a per-proof quantity, the count $n_t$ of
Definition~\ref{def:instance}; the prover pays only for those. Write
$m_t = \nu_t^{\max} + \kappa_t$; the type's \emph{capacity area} is
$s_t = 2^{m_t}$. The registry orders types by non-increasing capacity area
$s_1 \ge s_2 \ge \cdots \ge s_T$.
\end{definition}

Capacities are per-type rather than global because the generality is free
(slot areas differ across types anyway, since widths do) and workloads are
asymmetric: the address-space size $\sum_t s_t$ --- which sets the static
round count --- is dominated by the largest slots, so a type expected to
appear rarely should not be granted the row capacity of the workhorse
types. A uniform row capacity $\nu_t^{\max} \equiv \nu^{\max}$ is the
obvious special case; another natural one is the \emph{area-balanced}
registry, $\nu_t^{\max} + \kappa_t$ constant, where wider hashes get
proportionally fewer maximum rows and every slot has equal area.

Two clarifications on Definition~\ref{def:registry}. First, capacities are
powers of two \emph{by fiat}: slots must be subcubes of the global cube
(Lemma~\ref{lem:align} and everything built on it), and subcubes have
power-of-two size. This costs nothing --- capacity is virtual, so a true
maximum of, say, $3 \cdot 2^{10}$ is simply rounded up to $2^{12}$; the
never-used rows are skipped by the prover, and all genuine arbitrariness
lives in the counts $n_t$, which are unrestricted. Second, the registry
ordering is by capacity \emph{area} $s_t$ (rows $\times$ width), not by row
capacity alone --- a wider type with fewer maximum rows may sort first ---
and the ordering is itself only a device: it is the canonical way to obtain
the alignment property of Lemma~\ref{lem:align}, and in the area-balanced
case it is immaterial.

\begin{definition}[Instance]\label{def:instance}
A proof \emph{instance} consists of public \emph{counts}
$(n_t)_{t=1..T}$ with $0 \le n_t \le 2^{\nu_t^{\max}}$ --- arbitrary
integers, \emph{not} necessarily powers of two --- together with, per type
$t$, the witness table $Z_t \in \Ftwo^{n_t \times k_t}$: row $r < n_t$ holds
the full trace of invocation $r$ and independently satisfies
$(A_t z) \circ (B_t z) = z$. Rows $[n_t, 2^{\nu_t^{\max}})$ are \emph{dummy
rows}, identically zero.
\end{definition}

Independence of invocations means there are no constraints connecting
different rows, neither within a table nor across tables. The multi-table
layer attaches no statement to individual rows --- no input/output binding
(\S\ref{sec:statement}); in particular, dummy rows are referenced by
nothing at all.

If a workload exceeds a capacity, the proof is split (or a larger registry
tier is used, \S\ref{sec:statement}); capacities are a registry-design
question, not a protocol limit (\S\ref{sec:open}).

% ===========================================================================
\section{The union instance}\label{sec:union}

We compose the tables into one index space and one R1CS. The layout is
computed \textbf{once, at registry-definition time}, from the capacity
areas --- the per-proof counts never move anything.

\subsection{Static layout and alignment}

\paragraph{Layout convention (BatchMajor).} Within a slot, addresses use
the \emph{implemented} \code{WitnessLayout::BatchMajor} order
(\code{r1cs.rs}): trace variables are grouped into chunks of $128$; one
$128$-bit word holds one invocation's chunk; the words of a chunk are
contiguous across invocations. The address split is therefore
\emph{interleaved}:
\[
  \text{addr} = [\,\underbrace{7 \text{ in-word bits}}_{\text{column,
  low}} \mid \underbrace{\nu_t^{\max} \text{ row bits}}_{\text{invocation}}
  \mid \underbrace{\kappa_t - 7 \text{ chunk bits}}_{\text{column,
  high}}\,],
\]
i.e.\ slot $t$'s \emph{row coordinates} are $R_t = [7,\, 7{+}\nu_t^{\max})$
and its \emph{column coordinates} are the non-contiguous
$C_t = [0,7) \cup [7{+}\nu_t^{\max},\, m_t)$. The interleaving is
harmless: all the math below depends only on the \emph{partition} of a
slot's coordinates into (row, column, prefix) groups, never on their
contiguity. The convention earns its keep three ways: the univariate skip
runs over the low six in-word (column) bits \emph{exactly as today} ---
the existing skip kernels are byte-level layout-agnostic and run unchanged
(\S\ref{sec:skip}); chunk-columns are contiguous across invocations, the
stacking order the jagged commitment needs (\S\ref{sec:commit}), at
$128$-variable-chunk granularity; and the witness drivers already emit
this layout (with per-block padding coalesced), producing $z$, $Az$, and
$Bz$ during generation --- no matrix application on the hot path. Requires
$\kappa_t \ge 7$ (all current $\kappa_t \ge 14$).

Let $M$ be minimal with $2^{M} \ge \sum_t s_t$ (choose capacities so that
$\sum_t s_t = 2^{M}$ exactly if desired; a shortfall is harmless by
Remark~\ref{rem:padding}). Assign type $t$ the \emph{slot}
$[o_t, o_t + s_t)$ of $\{0,1\}^{M}$, where $o_t = \sum_{t' < t} s_{t'}$.

\begin{lemma}[Alignment]\label{lem:align}
With the descending-area ordering, $o_t$ is a multiple of $s_t = 2^{m_t}$.
Consequently, for a global index $x$ in slot $t$:
\[
  x|_{R_t} = \text{row index},\quad
  x|_{C_t} = \text{column index},\quad
  x[m_t{:}M] = p_t \;(\text{a fixed bit pattern}, \; p_t = o_t / 2^{m_t}),
\]
where $x|_S$ denotes the restriction of $x$ to the coordinate set $S$ and
$(R_t, C_t)$ is the interleaved row/column partition of the BatchMajor
convention above.
\end{lemma}
\begin{proof}
$o_t$ is a sum of terms $2^{m_{t'}}$ with $m_{t'} \ge m_t$, each divisible by
$2^{m_t}$.
\end{proof}

We call $p_t$ the \emph{slot prefix}: the Boolean string of length
$M - m_t$ that the top bits of \emph{every} address in slot $t$ are frozen
to (numerically $o_t / 2^{m_t}$, an exact division by the lemma). Note its
length varies per slot --- larger slots have shorter prefixes --- so
expressions like $\eq(\tau[m_t{:}M], p_t)$ always pair it with a
coordinate range of matching length. Example ($M = 6$, slot sizes
$32, 16, 16$): slots $[0,32)$, $[32,48)$, $[48,64)$ have prefixes
$p_1 = \mathtt{0}$, $p_2 = \mathtt{10}$, $p_3 = \mathtt{11}$.

So each slot is a subcube of $\{0,1\}^M$, selected by fixing the top
$M - m_t$ bits to the static prefix $p_t$, and \emph{within} the
subcube the index has the same (column, row) split as a single-table
instance. Distinct slots are disjoint subcubes. (The alignment property is
what all the math below needs; sorting is merely how the registry obtains
it once and for all.)

\subsection{The block-diagonal R1CS}

The union instance has constraint matrices
\[
  A = \mathrm{diag}\bigl(I_{2^{\nu_1^{\max}}} \otimes A_1,\; \dots,\;
                         I_{2^{\nu_T^{\max}}} \otimes A_T\bigr),
  \qquad B \text{ likewise}, \qquad C = I_{2^M},
\]
and witness $z \in \Ftwo^{2^M}$ = the concatenation of the per-type
witnesses in slot order (each table internally BatchMajor, per the
convention of \S\ref{sec:union}), with dummy rows zero. The satisfied relation is, as before, pointwise:
$(Az) \circ (Bz) = z$ over all of $\{0,1\}^M$.

\begin{remark}[Dummy rows are complete; padding is self-enforcing]
\label{rem:padding}
The R1CS is homogeneous --- constants enter only through the constant-pin
column --- so the all-zero row satisfies
$(A_t \cdot 0)\circ(B_t \cdot 0) = 0 = C_t \cdot 0$, with the constant pin
at $0$ marking the row as dummy; no witness-generation work is performed
for dummy rows. Soundness needs nothing from them either: nothing in the
system references dummy rows (the core binds no per-row IO at all,
\S\ref{sec:statement}), and under the jagged commitment
(\S\ref{sec:commit}) the region beyond height $n_t$ is zero \emph{by
construction} of the transport --- the prover has no freedom there to begin
with. In any gap region beyond the last slot, $A$ and $B$ have zero rows
and $C = I$, so the constraint reads $0 \cdot 0 = z(x)$, \emph{forcing}
zero. Padding therefore needs no soundness caveats; its only cost is
prover work, which \S\ref{sec:zc}--\ref{sec:lc} reduce to zero.
\end{remark}

There are \textbf{no selector polynomials} anywhere: disjointness of the
slots is the selection mechanism. The ``selector'' will surface only as
scalar $\eq$-prefix weights in the verifier's algebra, essentially for free.

% ===========================================================================
\section{Zerocheck over the union}\label{sec:zc}

Run \emph{exactly the existing zerocheck} (\S\ref{sec:pioponetable}) on the
union instance: sample $\tau \in \F^{M}$, prove
\begin{equation}\label{eq:zc-union}
  \sum_{x \in \{0,1\}^{M}} \eq(\tau, x)
     \bigl(\hat a(x)\hat b(x) - \hat c(x)\bigr) = 0 ,
\end{equation}
by an $M$-round sumcheck ($1$ univariate-skip round $+$ $M - 6$ ordinary
rounds), where now $a = Az$ etc.\ for the block-diagonal $A$. Note $M$ is
registry-determined, hence \emph{static}: a few more rounds than a
perfectly-sized instance, in exchange for a count-independent verifier.

\subsection{Why this is implicit batching}

Split the sum in \eqref{eq:zc-union} by slot. For $x$ in slot $t$,
$\eq(\tau, x)$ factors across the bit split of Lemma~\ref{lem:align}:
\[
  \eq(\tau, x) \;=\;
  \underbrace{\eq\bigl(\tau[m_t{:}M],\, p_t\bigr)}_{=: \, w_t \;\in\; \F}
  \;\cdot\;
  \eq\bigl(\tau[0{:}m_t],\, x[0{:}m_t]\bigr),
\]
and the constraint polynomial restricted to slot $t$ is exactly type
$t$'s own $\hat a_t \hat b_t - \hat c_t$. Hence
\begin{equation}\label{eq:implicit-batch}
  \eqref{eq:zc-union}
  \iff
  \sum_{t=1}^{T} w_t \cdot
  \underbrace{\sum_{x' \in \{0,1\}^{m_t}}
     \eq\bigl(\tau[0{:}m_t], x'\bigr)
     \bigl(\hat a_t \hat b_t - \hat c_t\bigr)(x')}_{\text{type } t\text{'s
     single-table zerocheck sum}} \;=\; 0 .
\end{equation}
That is: the union zerocheck \emph{is} a batched per-table zerocheck, with
batching coefficients $w_t = \eq(\tau[m_t{:}M], p_t)$ derived from the shared
randomness $\tau$ instead of independently sampled. There is no separate
batching protocol, no batching soundness lemma (soundness is the ordinary
single-instance zerocheck argument over $\{0,1\}^M$: if any invocation of any
table is wrong, the constraint polynomial is nonzero somewhere on the cube,
and \eqref{eq:zc-union} fails except with probability
$O(M/|\F|)$ over $\tau$ plus the sumcheck error), and no scheduling
machinery. Different tables simply live at different (static) addresses.

\subsection{The prover pays for declared work only}\label{sec:zc-sparse}

On every region where the witness --- hence $a$, $b$, and $c$ --- is
identically zero (dummy rows, useless columns, the gap), the summand of
\eqref{eq:zc-union} vanishes identically and contributes nothing to any
round polynomial. The \code{PaddingSpec} generalizes from today's single
$(\kappa, u)$ pair to a \textbf{run-list derived from the public counts}:
per slot, columns $[u_t, k_t)$ and rows $[n_t, 2^{\nu_t^{\max}})$ are dead
(this generalizes \code{useful\_bits} skipping from columns to rows). With
run-list--driven iteration the prover's sumcheck work is proportional to
$\sum_t n_t k_t$ --- the \emph{declared} area --- while the round count
stays the static $M$.

\subsection{The univariate skip survives globally --- unchanged}
\label{sec:skip}

The skip merges the six low variables of the sumcheck into one univariate
round; its prover win --- the six most expensive rounds computed directly
on $\Ftwo$-valued data --- requires only that each merged group of $64$
domain points be memory-contiguous. Under the BatchMajor convention the
global low six bits of any index are in-word \emph{column} bits for every
slot (all $\kappa_t \ge 7$, and every slot offset is a multiple of
$2^{m_t}$), exactly as in today's row-major layout: a skip group is $64$
adjacent trace variables of one invocation, one $u64$ of a $128$-bit
word. The existing round-1/round-2 skip kernels are byte-level
layout-agnostic and run \emph{unchanged} --- no new kernels, no
re-orientation. The zerocheck constraint is pointwise, so mixing slots
changes nothing else: one skip round, all slots at once, no per-table
scheduling. The skip's univariate coordinate sits among the \emph{column}
coordinates of the output point, precisely as today, so the lincheck's
quirky-point handling (\S\ref{sec:lc}) and the existing claim-point
conventions carry over; the coordinate is then eliminated by the existing
ring-switching stage before the jagged transport ever sees a point
(\S\ref{sec:commit}).

\subsection{What changes in code}

Only bookkeeping:
\begin{itemize}[nosep]
  \item \code{BlockR1cs} generalizes to the registry schedule of slots
        $\{(\text{base block}, \kappa_t, u_t, \nu_t^{\max})\}$ (one entry
        $\equiv$ today's struct); \code{apply\_a} applies each base block
        over its slot's declared rows.
  \item \code{PaddingSpec} $(\kappa, u)$ becomes the count-derived run-list
        of \S\ref{sec:zc-sparse}.
  \item The zerocheck sumcheck kernels: \textbf{unchanged, byte-identical}
        (pointwise and layout-agnostic; the code map confirmed the uniform
        \code{PaddingSpec} was the only single-table assumption, and the
        BatchMajor prover paths for BLAKE3/SHA-256 already exercise these
        kernels today).
\end{itemize}

\paragraph{Output.} A point $r \in \F^{M}$ and three claims
$\hat a(r), \hat b(r), \hat c(r)$, with $\hat c(r) = \hat z(r)$ a direct
witness claim.

% ===========================================================================
\section{Lincheck over the union}\label{sec:lc}

This is where the real (and modest) protocol work lives. The zerocheck left
two matrix claims, $\hat a(r)$ and $\hat b(r)$. For clarity we develop the
lincheck for the $A$-claim alone, $v_A := \hat a(r)$;
\S\ref{sec:batchAB} then observes that the $B$-claim is word-for-word
identical and that the two sumchecks merge into one at the cost of a single
random scalar.

\subsection{Closed form of the union matrix MLE}

Combining the block-diagonal structure with
\eqref{eq:factor-single} per slot and Lemma~\ref{lem:align}:
\begin{equation}\label{eq:factor-union}
  \hat A(r, c) \;=\; \sum_{t=1}^{T}
     \eq\bigl(r[m_t{:}M], p_t\bigr)\,
     \eq\bigl(c[m_t{:}M], p_t\bigr)\,
     \eq\bigl(r|_{R_t},\, c|_{R_t}\bigr)\,
     \hat A_t\bigl(r|_{C_t},\, c|_{C_t}\bigr),
\end{equation}
\paragraph{Derivation of \eqref{eq:factor-union}.} By definition
$\hat A(r,c) = \sum_{x,y \in \{0,1\}^M} \eq(r,x)\,\eq(c,y)\,A[x,y]$, and
the MLE is linear in the matrix, so split $A$ by slot; cross-slot entries
vanish, and slot $t$'s part is supported on pairs with both prefixes frozen
at $p_t$. Since $\eq$ is a product over coordinates, it factors over any
partition of them: for $x$ in slot $t$,
$\eq(r,x) = \eq(r[m_t{:}M], p_t)\cdot\eq(r|_{R_t}, x|_{R_t})
\cdot\eq(r|_{C_t}, x|_{C_t})$ --- $\eq$ factors over \emph{any} partition
of the coordinates, contiguous or not --- and the prefix factor ---
constant over the slot --- pulls out of the sum (likewise for $c$): these
are the first two factors of \eqref{eq:factor-union}, the multilinear
indicators ``$r$ addresses slot $t$'' and ``$c$ addresses slot $t$''. What
remains is the MLE of $I \otimes A_t$ in the within-slot coordinates, and
the MLE of a Kronecker product is the product of the factors' MLEs: the
identity contributes $\eq(r|_{R_t}, c|_{R_t})$
(``same invocation'') via
\begin{equation*}
  \sum_{w \in \{0,1\}^{\nu}} \eq(a, w)\,\eq(b, w) \;=\; \eq(a, b),
\end{equation*}
and the base block contributes $\hat A_t$. Note each term partitions the
\emph{same} shared coordinates by its own $(R_t, C_t, \text{prefix})$ ---
legal, since every term is simply a function of all $M$ coordinates.
Sanity check: at a Boolean pair $(x, y)$, every term except the one for
the slot actually containing $x$ dies at its prefix factor, and the
surviving term reproduces the matrix entry.

The verifier can evaluate the scalar prefix weights in $O(T \cdot M)$
additions and multiplications; the per-type base-block evaluations cost
$O(\sum_t \nnz(A_t, B_t))$ via the existing \code{LincheckCircuit} comb
machinery --- still never touching any expanded matrix.

\subsection{Per-slot collapse}

Substituting \eqref{eq:factor-union} into
$v_A = \sum_{c \in \{0,1\}^M} \hat A(r, c)\, \hat z(c)$ and
collapsing, per slot, the sum over the row coordinates $c|_{R_t}$ into a
partial evaluation of that slot's witness MLE (exactly as in the
single-table case):
\begin{equation}\label{eq:lc-claim}
  v_A \;=\; \sum_{t=1}^{T} w_t(r) \sum_{c' \in \{0,1\}^{\kappa_t}}
      \mathrm{comb}^A_t(c')\; y_t(c'),
  \qquad
  \begin{aligned}
    w_t(r) &= \eq\bigl(r[m_t{:}M], p_t\bigr),\\
    y_t(c') &= \hat z_t\bigl(r|_{R_t},\, c'\bigr),
  \end{aligned}
\end{equation}
where
$\mathrm{comb}^A_t(c') = \sum_{x'} \eq(r|_{C_t}, x')\, A_t[x', c']$
is the base block folded by the weight table of the within-block
constraint coordinates of $r$ --- the \emph{quirky} table, since $C_t$
contains the univariate skip coordinate (\S\ref{sec:skip}), computed
through the type's \code{LincheckCircuit} \emph{exactly as in today's
single-table code} (the code's ``semantic \code{QuirkyPoint}'' is
layout-independent) --- and $\hat z_t(\cdot,\cdot)$ is the MLE of slot
$t$'s (zero-padded) witness, rows-coordinates first. The prover
materializes each $y_t$ by folding slot $t$'s \emph{declared rows} with
the standard $\eq$ table of the row coordinates $r|_{R_t}$ ---
$O(\sum_t n_t k_t)$ in total, the declared area, analogous to today's
single outer fold. The division of labor (quirky on the comb side,
standard $\eq$ on the row-fold side) is \emph{identical} to today's
single-table lincheck.

Note the weights $w_t(r)$ are \emph{fixed known scalars} (determined by $r$
and the registry), not fresh randomness: \eqref{eq:lc-claim} is a single
claim about a fixed weighted sum, so no batching randomness is needed
\emph{across slots} (contrast \S\ref{sec:batchAB}, where combining the
$A$- and $B$-claims genuinely does use a random $\alpha$).

\subsection{The multi-size inner sumcheck}\label{sec:multisize}

What remains is one sumcheck for \eqref{eq:lc-claim}: a sum of $T$
sub-sums over domains of different sizes $2^{\kappa_t}$,
$\kappa_t \in \{14, \dots, 17\}$. Two mechanisms, in order of preference:

\paragraph{(a) $\eq$-embedded late join (recommended).}
Let $K = \max_t \kappa_t$. Embed each sub-sum into $\{0,1\}^{K}$ using a
shared auxiliary point $\sigma \in \F^{K}$:

\begin{lemma}[$\eq$-embedding]\label{lem:embed}
For any $f : \{0,1\}^{\kappa} \to \F$, any $K \ge \kappa$, and any
$\sigma \in \F^{K}$:
\[
  \sum_{c \in \{0,1\}^{K}}
     \eq\bigl(\sigma[\kappa{:}K],\, c[\kappa{:}K]\bigr)\; f\bigl(c[0{:}\kappa]\bigr)
  \;=\; \sum_{c' \in \{0,1\}^{\kappa}} f(c'),
\]
by \eqref{eq:eqsum} applied to the top $K - \kappa$ coordinates.
\end{lemma}

\begin{remark}[Characteristic 2 forces this]\label{rem:char2}
The textbook alternative --- replicate the small summand over the top
coordinates and rescale by $2^{-(K-\kappa)}$ --- is unavailable: in
$\F_{2^{128}}$, $2 = 0$. The $\eq$-weighting of Lemma~\ref{lem:embed} is the
canonical char-2-safe embedding, and the same remark applies anywhere else in
Flock where sums over different-sized cubes are ever combined.
\end{remark}

The sumcheck then has $K$ rounds, binding coordinates high to low. In a round
where coordinate $j \ge \kappa_t$ is bound, type $t$'s term contributes
$w_t \cdot \bigl(\prod_{i>j} \eq(\sigma_i, \rho_i)\bigr)\,\eq(\sigma_j, X)\,
 v_t$ --- a \emph{known scalar} times a linear polynomial, where
$v_t = \sum_{c'} \mathrm{comb}^A_t(c')\,y_t(c')$ is the sub-claim value
(which the prover sends alongside the round data; the verifier checks
$\sum_t w_t v_t = v_A$). Once coordinate $\kappa_t - 1$ is reached, type $t$
``joins'' and contributes its ordinary product-sumcheck rounds over
$(\mathrm{comb}^A_t, y_t)$, scaled by its accumulated prefix. With
$\kappa \in \{14..17\}$ all joins happen within the first
$K - \min_t \kappa_t \le 3$ rounds. At the end, the verifier holds, per
type, $\mathrm{comb}^A_t$ evaluated at the final point (computed itself via
the comb machinery at $O(\nnz)$ per type) and an evaluation claim
$y_t(c^*_t) = \hat z_t(r|_{R_t}, c^*_t)$ --- one witness claim
per type.

\paragraph{(b) Fallback: full-domain lincheck.} Alternatively run the
lincheck over all $M$ column bits, with summand $\hat A(r,c)\hat z(c)$ via
\eqref{eq:factor-union} --- zero new sumcheck machinery, at the cost of one
extra pass to materialize $\hat A(r, \cdot)$ on the (declared part of the)
cube. We spec (a) and keep (b) as the de-risking fallback; the delta between
them is a benchmark, not an architecture question.

\subsection{The \texorpdfstring{$B$}{B}-claim, and batching the two
sumchecks}\label{sec:batchAB}

Everything above applies verbatim to the $B$-claim $v_B := \hat b(r)$, with
$\mathrm{comb}^B_t$ folding $B_t$ instead of $A_t$. Crucially, the folded
witness vectors $y_t$ \emph{do not change} --- they depend only on $r$ and
the witness, never on the matrix. The two multi-size sumchecks therefore
range over identical domains against identical $y_t$'s and differ only in
the comb factor, so they merge into one: squeeze a fresh $\alpha$ from the
transcript (after both claims are bound) and run the sumcheck of
\S\ref{sec:multisize} once for the combined claim
\[
  \alpha\, v_A + v_B \;=\; \sum_{t=1}^{T} w_t(r) \sum_{c'}
    \bigl(\alpha\,\mathrm{comb}^A_t + \mathrm{comb}^B_t\bigr)(c')\;
    y_t(c') .
\]
This is standard random-linear-combination batching of two claims (error
$1/|\F|$), costs one scalar and zero extra rounds, and the combined comb
$\alpha\,\mathrm{comb}^A_t + \mathrm{comb}^B_t$ is exactly what the
existing \code{fold\_alpha\_batched} computes --- the code's
$\alpha$-comb is this batching, seen from the single-table special case.
The end-of-sumcheck outputs --- per-type comb evaluations and the witness
claims $y_t(c^*_t)$ --- are shared between the two claims, so batching
adds nothing there either.

\subsection{Verifier cost summary}

Rounds: $M$ (zerocheck) $+$ $K \le 17$ (lincheck) --- both
registry-static, independent of $T$ and of the counts. Scalars: $3$
zerocheck claims $+$ $T$ sub-claim values $+$ $T$ witness claims. Matrix
work: $O(\sum_t \nnz(A_t, B_t))$, linear in the number of types --- the one
term that grows if the registry grows large; see \S\ref{sec:open} for the
delegation escape hatch. Prefix weights: $O(T \cdot M)$. Weight
\emph{structure}, join schedule, comb structures, and the registry digest
are all precomputable; the counts enter only through run lists and fold
extents. One bookkeeping note: a union-level claim such as
$\hat c(r) = \hat z(r)$ decomposes via the prefix weights into per-slot
claims $\hat z_t(\cdot)$, which join the same transport batch as everything
else.

% ===========================================================================
\section{The commitment layer: jagged transport}\label{sec:commit}

The PIOP above emits evaluation claims on the per-slot virtual polynomials
$\hat z_t$ (from the lincheck and from the decomposed union claims; a
future lookup/bus layer will add its own, \S\ref{sec:open}). The
commitment layer
turns these into openings of \emph{one} committed polynomial while
committing \textbf{only the actual data}: the capacity-sized address space
of \S\ref{sec:union} exists purely virtually.

\paragraph{Pipeline position: ring-switch first, then jagged.} Claims
emitted by the PIOP are bit-level and carry the univariate skip
coordinate among their low (in-word) coordinates. They first pass through
the \emph{existing} ring-switching stage exactly as today
(\S\ref{sec:pioponetable}): ring-switching eliminates the seven in-word
coordinates --- the skip coordinate included --- yielding per-slot
multilinear claims on the \emph{packed} polynomial, whose variables are
the batch and chunk coordinates. Ring-switching is per-claim local, so it
needs no union-awareness. Its prover folds are, however, subject to the
same declared-work discipline as everything else: the $\hat s_v$ capture
for the zerocheck claim is fused into the zerocheck pass and inherits its
run-lists (an explicit checkpoint when generalizing that pass), and the
folds for the per-slot lincheck claims touch only their own slot's
declared rows, $O(n_t k_t)$ each. The jagged transport therefore operates purely
on packed data at ordinary multilinear points --- matching both the
\code{jagged-assist} API (plain multilinear
$(z_{\mathrm{row}}, z_{\mathrm{col}}, z_{\mathrm{index}})$ slices) and
its flagship bench ($2^{30}$-bit \emph{packed} witness). At the packed
level the jagged grid reads: columns $=$ chunk-columns, rows $=$
invocations, heights $= n_t$ \emph{words}.

\paragraph{The jagged scheme.} Stack only the \emph{useful} columns of the
\emph{declared} rows of all tables contiguously (column-major per table ---
the \code{BatchMajor} witness layout, added in anticipation of exactly
this): the committed size is $\sum_t n_t \lceil u_t/128\rceil \cdot 128$ bits
(stacking granularity is the $128$-variable chunk-column, so useless
columns are dropped per chunk; the sub-chunk remainder is negligible),
with all chunk-columns of slot $t$ having height exactly $n_t$ ---
non-power-of-two counts are native. Because the union's virtual layout is
itself BatchMajor (\S\ref{sec:union}), the committed stack is simply the
union buffer with dummy rows and useless chunks omitted --- commitment
and address space share one orientation, with no transpose anywhere in
the pipeline. Commit
the dense stack with Ligerito, and transport each virtual claim
$\hat z_t(\rho)$ to the dense polynomial via the jagged sumcheck
(\code{pcs/jagged.rs}, branch \code{jagged-assist}):
\code{col\_prefix\_sums} --- derived per proof from the public counts ---
encodes the layout, and the branching-program evaluator $\hat f_t$, with
the ``assist'' delegation, keeps the verifier cheap. The virtual table
beyond height $n_t$ is zero \emph{by definition} of this transport, which
is what discharged the soundness obligation of Remark~\ref{rem:padding}.
All transported claims are batched (random linear combination) into one
jagged sumcheck and one Ligerito batch opening; the PIOP is oblivious to
all of it.

One consequence to watch: with per-proof counts, the jagged verifier's
branching-program data is per-proof rather than fully precomputed --- this
is precisely where the assist delegation earns its keep; the
registry-determined part of the structure remains precomputable.

\begin{remark}[The direct-commitment harness]\label{rem:harness}
At full utilization ($n_t = 2^{\nu_t^{\max}}$ for all $t$) the union buffer
is itself a dense power-of-two array, and it can be committed directly with
today's \code{pcs::commit}; per-slot claims then transport \emph{trivially}:
for $\rho \in \F^{m_t}$, let $x^\star \in \F^{M}$ be $\rho$ in coordinates
$[0, m_t)$ and the bits of $p_t$ in coordinates $[m_t, M)$; then
$\hat z(x^\star) = \hat z_t(\rho)$. (Proof: decompose $\hat z$ by slot; slot
$t$'s prefix factor is $\eq(p_t, p_t) = 1$; for $t' \ne t$ with
$m_{t'} \le m_t$ the factor is $\eq$ of two distinct Boolean strings $= 0$;
for $m_{t'} > m_t$, if the top $M - m_{t'}$ bits of $p_t$ equaled $p_{t'}$,
slot $t$'s address range would lie inside slot $t'$'s, contradicting
disjointness.) This path is \textbf{not a protocol mode} and does not
appear in the wire format. It exists as (i) the differential-testing
harness for the union PIOP before and after the jagged layer lands, and
(ii) a cheaply resurrectable fallback should the jagged composition stall
--- the PIOP is identical under both transports.
\end{remark}

% ===========================================================================
\section{Statement, transcript, wire format}\label{sec:statement}

\paragraph{Statement digest and transcript.} The \emph{registry digest} ---
format version, the ordered type list $(\text{type}_t, \kappa_t, u_t,
\nu_t^{\max})$, and each base block's existing digest --- is one precomputed
32-byte value. Per proof, the transcript absorbs, before any challenge is
squeezed: the registry digest, the counts vector $(n_t)_t$, and the
commitment root. New domain-separation labels (\code{flock-mixed-v1}, plus
labels for the multi-size lincheck and the jagged transport) join the
existing family. Counts affect only scalar bookkeeping --- run lists, fold
extents, \code{col\_prefix\_sums} --- never the verifier's control flow.

\paragraph{What a multi-table proof attests.} Exactly this: \emph{the
commitment opens to tables with the declared counts, every declared row of
which satisfies its type's hash relation.} The core deliberately carries
no statement about the hashed values themselves --- no per-invocation
input/output binding. All dataflow --- between tables, and between tables
and the public statement --- is the job of a planned \emph{lookup/bus}
layer built on top of this core (\S\ref{sec:open}): the core's only
obligations toward it, both already met, are fixed column offsets per type
(the \code{ChainLayout} region data) and cheap per-slot claims into the
transport batch. Until the bus lands, mixed proofs are infrastructure and
benchmarks, not application statements; today's single-type statements
(hash chains, Merkle paths) are untouched and keep their existing
claim-level glue.

\paragraph{Capacity tiers.} If workloads vary wildly, the registry may
offer a few \emph{capacity tiers} (small/medium/large instantiations of the
same type list); selecting a tier selects a registry digest. Each doubling
of capacity costs one static sumcheck round; see \S\ref{sec:open}.

\paragraph{Wire format.} \code{proof\_io}: bump \code{VERSION} to 5; replace
the single \code{hash\_kind} in the bundle with the registry (tier) id and
the counts vector. CLI: \code{flock\_chain} gains
a mix argument (e.g.\ \code{--mix blake3=1000,sha2=4096}); the existing
single-hash flags become single-type mixes internally.

% ===========================================================================
\section{Cost accounting}\label{sec:costs}

Let $N_{\mathrm{decl}} = \sum_t n_t k_t$ (declared area),
$N_{\mathrm{com}} = \sum_t n_t u_t$ (committed size), $T$ = types.

\begin{center}
\begin{tabular}{@{}ll@{}}
\toprule
 & \textbf{cost} \\
\midrule
Witness generation & per-type closures over declared rows; nothing for dummies \\
Commit & Ligerito on $N_{\mathrm{com}}$ (actual data, useful columns only) \\
Zerocheck prover & $O(N_{\mathrm{decl}})$; static $M$ rounds; skip intact \\
Lincheck prover & $O(N_{\mathrm{decl}})$ folds $+$ $K$-round multi-size SC \\
Ring-switch folds & $O(N_{\mathrm{decl}})$ (fused into the zerocheck pass; per-slot for \LC claims) \\
Jagged transport prover & $\tilde O(N_{\mathrm{com}})$ (one batched transport SC) \\
Opening & one batched Ligerito open \\
\midrule
Verifier rounds & $M + K$ $+$ transport --- registry-static, count-independent \\
Verifier matrix work & $O(\sum_t \nnz(A_t,B_t))$ \\
Verifier bookkeeping & $O(T \cdot M)$ weights $+$ per-proof jagged data (assisted) \\
Proof size & $+\,O(T)$ scalars over a single-table proof of equal area \\
\bottomrule
\end{tabular}
\end{center}

The headline: \textbf{the prover pays for declared work only} --- every
per-invocation cost (witness gen, zerocheck folding, commit) matches the
single-table system on the same workload; capacity costs only
$M - \log_2 N_{\mathrm{decl}}$ extra sumcheck rounds, and mixing costs
$O(T)$ scalars plus the (new, but size-independent-per-invocation) jagged
transport. Relative to today's proofs, the committed area actually
\emph{shrinks} by the useful-column ratio ($\approx 4$--$6\%$).

% ===========================================================================
\section{Roadmap}\label{sec:roadmap}

Jagged is on the critical path for everything, so it lands first, in the
smallest possible context.

\begin{description}[leftmargin=1.6em, style=nextline]
  \item[Phase 0 --- enablers (no behavior change).]
    Registry/schedule types; \code{BlockR1cs} $\to$ slot schedule
    (single-table $=$ one slot; all existing tests must pass unchanged);
    run-list \code{PaddingSpec}; registry digest. Gate: existing
    single-hash proofs byte-stable (modulo version field) and benchmarks
    flat.
  \item[Phase 1 --- jagged under the existing single-table system.]
    Wire the \code{jagged-assist} transport beneath today's statements
    (one-type registry, the existing BatchMajor prover paths): same
    proofs, new commitment path, so every existing test and benchmark is
    a correctness and throughput oracle. The zerocheck/lincheck kernels
    need no changes (\S\ref{sec:skip}); the work is the composition
    layer: slotting the jagged transport between the existing ring-switch
    stage and the Ligerito opening, and multi-claim RLC batching into one
    transport sumcheck (\S\ref{sec:open}). Gate: single-hash throughput
    within an agreed margin of the direct path; chunk-granular area
    saving realized.
  \item[Phase 2 --- the union PIOP and dynamic counts (mixed proofs ship).]
    Union zerocheck (count-derived run lists); multi-size lincheck
    (\S\ref{sec:multisize}(a), with (b) as fallback); per-slot witness
    closures in the existing drivers; counts in the
    transcript; wire format v5; CLI \code{--mix}. Internal milestone: the
    union PIOP differentially tested at full utilization against the
    direct-commitment harness (Remark~\ref{rem:harness}). Gates: mixed
    BLAKE3$+$SHA-256 proof verifies end-to-end; prover time tracks
    \emph{declared} work across utilization levels; verifier time flat
    across utilization levels; throughput within noise of the weighted
    average of the single-hash provers on the same workload.
\end{description}

Independent invocations keep every phase per-table testable: any mixed
proof can be differentially tested against the corresponding single-table
proofs on the same witnesses.

% ===========================================================================
\section{Open questions}\label{sec:open}

\begin{enumerate}[nosep]
  \item \textbf{Jagged claim batching (Phase 1).} The jagged kernel's
        public API (branch \code{jagged-assist}) transports exactly one
        sparse claim per invocation. Because ring-switching runs first
        and eliminates the in-word coordinates --- univariate skip
        coordinate included (\S\ref{sec:commit}) --- the transport only
        ever sees ordinary multilinear points over packed-word variables,
        so \emph{no point-format adapter is needed}. What the composition
        layer must add is claim \emph{batching} --- a cost matter, not a
        soundness one. A mixed proof emits $\approx 2T$ claims at
        distinct points ($T$ from decomposing the zerocheck's
        $\hat c(r) = \hat z(r)$ claim by prefix weights, $T$ from the
        per-slot lincheck outputs, plus glue on single-type statements),
        and the transport's dense-domain sumcheck and assist are paid per
        claim; RLC-combining the per-claim $\eq$-weight tables under
        transcript coefficients collapses this to one transport sumcheck
        --- the same multi-claim front-end \code{open\_batch\_mixed}
        already provides at the Ligerito layer. A per-claim transport v0
        is sound and an acceptable stepping stone inside Phase 1;
        batching follows as the optimization.
  \item \textbf{Capacity sizing.} How many tiers, and how large? Each
        doubling of capacity costs one static sumcheck round and slightly
        larger bookkeeping tables; exceeding capacity forces a proof
        split. Leaning: a single generous area-balanced tier
        (\S\ref{sec:model}), adding tiers only on demonstrated demand;
        confirm against a workload model at Phase 2.
  \item \textbf{Folding over empty capacity --- an implementation
        pattern, not a hot loop.} The union fold factorizes per slot:
        binding variables low to high, slot $t$ folds its own dense,
        declared data for its first $m_t$ rounds (today's kernels,
        verbatim) and then collapses to a single scalar; in every later
        round it contributes $O(1)$ scalar work (the accumulated value
        times the appropriate $\eq$ factors --- the implicit-batching
        identity \eqref{eq:implicit-batch}, operationally). So ``prover
        pays declared work'' is delivered by architecture --- per-slot
        dense sub-folds plus an $O(T)$-per-round scalar tail --- with no
        sparse iteration inside any hot loop. The same discipline covers
        the ring-switching folds (\S\ref{sec:commit}): the fused
        $\hat s_v$ capture and the per-slot claim folds must skip dummy
        rows via the same run-lists. Residual work: the round-polynomial
        accounting at slot-collapse boundaries, and an adversarial
        benchmark at low utilization.
  \item \textbf{Many types.} Verifier matrix work is linear in registry
        size ($\sum_t \nnz$). If the registry grows past a handful of
        types, the clean fix is Spartan-style preprocessing: commit each
        base block's MLE once at registry-definition time and replace the
        verifier's comb evaluation with one PCS opening per type;
        assist-style delegation is the lighter-weight alternative. Not
        needed initially.
  \item \textbf{The lookup/bus layer (the designated follow-on).} All
        dataflow --- table-to-table communication \emph{and} binding
        hashed values to the public statement --- is to be provided by a
        lookup/bus argument over the committed columns: rows emit and
        consume tuples (e.g.\ (type, input, output)); a multiset-balance
        argument (grand-product or log-derivative/LogUp style, realized as
        additional sumchecks over committed columns) enforces that every
        consumed tuple was emitted; public values enter as statement-side
        injections into the bus. This integrates at claim level --- the
        bus sumchecks emit evaluation claims into the same transport batch
        --- so the multi-table core needs nothing from it beyond what
        already holds (fixed column offsets per type, cheap per-slot
        claims). It changes no decision in this document and deserves its
        own design document.
  \item \textbf{Internal-hash generality.} Orthogonal axis (Merkle/FS hash
        is hard-coded SHA-256); out of scope here, noted for completeness.
\end{enumerate}

% ===========================================================================
\appendix
\section{Code touchpoint map}\label{app:code}

\begin{center}
\small
\begin{tabular}{@{}lll@{}}
\toprule
\textbf{Area} & \textbf{Location (today)} & \textbf{Change} \\
\midrule
Instance & \code{flock-core/src/r1cs.rs} (\code{BlockR1cs}) & slot schedule; \code{apply\_a} per slot \\
Zerocheck & \code{flock-core/src/zerocheck.rs} (\code{PaddingSpec}) & count-derived run lists; kernel unchanged \\
Lincheck & \code{flock-core/src/lincheck.rs} & multi-size inner SC; prefix weights; \\
 & (\code{LincheckCircuit}, \code{fold\_alpha\_batched}) & per-type comb reuse \\
Commitment & \code{flock-core/src/pcs/jagged.rs} $+$ branch & wire in as the claim-transport layer \\
Claim batching & \code{flock-core/src/pcs.rs} (\code{open\_batch\ldots}) & feeds the jagged transport batch \\
Test harness & \code{flock-core/src/pcs/commit.rs} & direct commit at full utilization \\
 & & (Remark~\ref{rem:harness}; tests only) \\
Statement & \code{flock-core/src/r1cs.rs} (\code{statement\_digest}) & registry digest $+$ counts \\
Witness gen & \code{flock-prover/src/r1cs\_hashes/common.rs} & drivers over per-slot closures \\
Statement glue & \code{flock-prover/src/r1cs\_hashes/chain\_common.rs} & unchanged (single-type statements); \\
 & & superseded later by the bus layer \\
Wire/CLI & \code{flock-prover/src/proof\_io.rs}, \code{bin/flock\_chain.rs} & v5; registry id $+$ counts; \code{--mix} \\
\bottomrule
\end{tabular}
\end{center}

\end{document}
